\chapter{Reconstruction Baseline and Codebook Usage}
\label{ch:rq1}

Before any perturbation is introduced, this chapter establishes the clean
ImageNet baseline for the two studied tokenizers. It measures how well each
tokenizer reconstructs unmodified validation images and which empirical token
distribution it produces. Both tokenizers have the same nominal vocabulary
size, $K=16{,}384$, and produce a $16\times16$ token grid from a
$256\times256$ image, making their outputs directly comparable.

\section{Experimental Setup}
\label{sec:rq1_experimental_setup}

For each image $x$, the tokenizer produces a reconstruction:

\begin{equation}
\hat{x} = D(Q(E(x))),
\end{equation}

where $E$ is the encoder, $Q$ is the vector-quantization step, and $D$ is the decoder. Reconstruction quality is measured using the metrics defined in Chapter~\ref{ch:datasets_models_metrics}: PSNR, SSIM, LPIPS, FID, sFID, precision, recall, and Inception Score. Codebook usage is computed from the discrete token grids over all 50,000 validation images. Since each image produces a $16 \times 16$ grid, each tokenizer contributes:

\begin{equation}
50{,}000 \times 16 \times 16 = 12{,}800{,}000
\end{equation}

token assignments for the ImageNet baseline.

The codebook analysis reports both raw and effective vocabulary use. The raw active count is the number of codebook entries assigned at least once. The effective active count uses a minimum-frequency threshold of $10^{-6}$ of all token assignments, corresponding to at least 13 assignments out of 12.8 million tokens in this experiment. Entropy and perplexity summarize the effective size of the empirical distribution, and top-$k$ mass measures how much of the token stream is captured by the most frequent codes.

The experiment uses the same ImageNet validation images and deterministic
inference for both models. Each exporter applies its own resize filter before
the $256\times256$ center crop; Appendix~\ref{app:computational_environment}
documents this preprocessing difference.

\section{Reconstruction Baseline}
\label{sec:rq1_reconstruction_baseline}

Table~\ref{tab:rq1_reconstruction_metrics} reports the ImageNet reconstruction metrics for LlamaGen and VQGAN. LlamaGen improves the paired-image metrics (PSNR, SSIM, LPIPS) and the reference-conditioned distribution metrics (FID, sFID, precision, and recall). Inception Score is retained as a diagnostic of classifier confidence and class diversity, not treated as a paired reconstruction metric or as having a universally preferred direction. The remaining metrics make LlamaGen the stronger clean reconstruction baseline on ImageNet.

\begin{table}[t]
    \centering
    \small
    \setlength{\tabcolsep}{4pt}
    \begin{tabular}{lrrrrrrrr}
        \textbf{Tokenizer} & \textbf{PSNR} & \textbf{SSIM} & \textbf{LPIPS} & \textbf{FID} & \textbf{sFID} & \textbf{Precision} & \textbf{Recall} & \textbf{IS} \\ \hline
        LlamaGen & 20.79 & 0.558 & 0.228 & 2.19 & 5.01 & 0.978 & 0.995 & 52.70 \\
        VQGAN    & 19.65 & 0.489 & 0.286 & 5.06 & 6.49 & 0.920 & 0.973 & 46.76 \\
    \end{tabular}
    \caption{ImageNet reconstruction metrics for the two tokenizers. Higher PSNR, SSIM, precision, and recall are preferred; lower LPIPS, FID, and sFID are preferred. Inception Score is reported as a diagnostic and is not assigned a preferred reconstruction direction.}
    \label{tab:rq1_reconstruction_metrics}
\end{table}

The PSNR gap is modest in absolute terms, but it is consistent with the other metrics. LlamaGen reaches a mean PSNR of 20.79 compared with 19.65 for VQGAN, and a mean SSIM of 0.558 compared with 0.489. LPIPS also favors LlamaGen: 0.228 compared with 0.286 for VQGAN. Because LPIPS is computed in a learned feature space, the reconstruction difference extends beyond pixel-level measurements. The distributional metrics agree with the same direction: LlamaGen obtains FID 2.19 and sFID 5.01, while VQGAN obtains FID 5.06 and sFID 6.49.

The corresponding per-image standard deviations are 3.18 dB, 0.167, and
0.058 for LlamaGen PSNR, SSIM, and LPIPS, and 3.41 dB, 0.176, and 0.066 for
VQGAN. These standard deviations describe variation across images, not
uncertainty in the dataset means or the fraction of images on which one model
wins. FID and related distribution metrics do
not decompose into per-image observations and are consequently reported as
dataset-level estimates.

The published LlamaGen comparison reports PSNR $20.79$, SSIM $0.675$, and
reconstruction FID $2.19$ for VQ-16, versus PSNR $20.00$, SSIM $0.629$, and
FID $4.99$ for VQGAN-f16-16384~\cite{sun2024llamagen}. Our LlamaGen PSNR and
FID agree to the displayed precision. The SSIM values are not directly
comparable. The public LlamaGen evaluation script sets the SSIM data range to
$2.0$ for images in $[0,1]$, whereas our exporters set it to $1.0$ and use
scikit-image's default window with the last axis as the color channel. Our
VQGAN PSNR and FID also differ from the published values;
the resize filters differ between exporters, so these numbers do not establish
the cause of that offset. LlamaGen's published $97\%$ usage is measured in a
65{,}536-token queue, whereas our $100\%$ denotes codes observed at least once
across 12.8 million assignments.

Figure~\ref{fig:rq1_reconstruction_metric_comparison} visualizes the same comparison across metrics that differ in scale and direction. Excluding Inception Score, whose interpretation is not paired to a reference image, the reconstruction metrics consistently favor LlamaGen.

\begin{figure}[t]
    \centering
    \includegraphics[width=0.98\textwidth]{ch04_fig_reconstruction_metric_comparison.png}
    \caption{ImageNet reconstruction metric comparison for VQGAN and LlamaGen. Each panel names whether higher or lower values are preferred; the comparison should be read by direction within each metric, not by comparing bar heights across panels.}
    \label{fig:rq1_reconstruction_metric_comparison}
\end{figure}

Figure~\ref{fig:rq1_reconstruction_triplets} shows how much detail and structure
each tokenizer preserves through the discrete bottleneck on three paired
examples.

\begin{figure}[t]
    \centering
    \includegraphics[width=0.9\textwidth]{ch04_fig_reconstruction_triplets.png}
    \caption{Qualitative ImageNet reconstruction examples. Each row shows the original image, the VQGAN reconstruction, and the LlamaGen reconstruction.}
    \label{fig:rq1_reconstruction_triplets}
\end{figure}

\section{Codebook Usage}
\label{sec:rq1_codebook_usage}

The metrics above show that LlamaGen reconstructs the ImageNet validation images more accurately than VQGAN under the recorded exporter settings. The following analysis examines whether the two tokenizers also differ in how they use their discrete vocabularies. Table~\ref{tab:rq1_codebook_usage} reports the main usage statistics.

\begin{table}[t]
    \centering
    \small
    \setlength{\tabcolsep}{4pt}
    \resizebox{\textwidth}{!}{%
    \begin{tabular}{lrrrrrrr}
        \textbf{Tokenizer} & \textbf{Codebook size} & \textbf{Active} & \textbf{Effective active} & \textbf{Perplexity} & \textbf{Top-10} & \textbf{Top-100} & \textbf{Top-500} \\ \hline
        LlamaGen & 16{,}384 & 16{,}384 (100.0\%) & 16{,}384 (100.0\%) & 16{,}179.19 & 0.003 & 0.012 & 0.046 \\
        VQGAN    & 16{,}384 & 972 (5.9\%)       & 971 (5.9\%)       & 910.46    & 0.022 & 0.165 & 0.665 \\
    \end{tabular}
    }
    \caption{ImageNet codebook usage summary. Active counts report how many nominal codes appear on ImageNet; effective active counts use the threshold of at least 13 assignments. Top-$k$ values report the fraction of all token assignments covered by the most frequent $k$ codes.}
    \label{tab:rq1_codebook_usage}
\end{table}

LlamaGen uses all 16,384 codebook entries on ImageNet with a perplexity of 16,179. This means that the empirical token distribution is close to uniform over the full nominal vocabulary. VQGAN uses only 972 of 16,384 codes, or about 5.9\% of the nominal codebook. Liu et al.~\cite{liu2024elucidating} previously reported approximately 1{,}000 used codes for this released VQGAN checkpoint on ImageNet training data; our validation-set count confirms that underuse under a different sample and protocol. Its perplexity of 910 is close to the raw active count, so the active subset is small but not dominated by only a handful of codes.

LlamaGen's top-500 codes account for only 4.6\% of all token assignments; VQGAN's top-500 account for 66.5\%. Thus, although the two tokenizers expose the same nominal vocabulary size, the sequences they would provide to a downstream autoregressive model have very different empirical vocabularies. The downstream model is trained on this empirical token stream, not on the unused portion of the nominal codebook.

Entropy gives an idealized coding interpretation of the same result. Since
$H=\log(\operatorname{PPL})$, the observed perplexities correspond to about
$9.69$ nats ($13.98$ bits) per token for LlamaGen and $6.81$ nats ($9.83$
bits) per token for VQGAN. For a 256-token grid, these are approximately
$3{,}579$ and $2{,}516$ idealized bits per image under an optimal code matched
to each empirical marginal distribution. LlamaGen's empirical marginal
entropy is therefore about 42\% higher. These marginal entropies describe ideal
code lengths for the observed token frequencies. They do not measure semantic
content or the conditional code length of a trained autoregressive model.

Figure~\ref{fig:rq1_code_usage_sorted} shows the full rank-frequency curve. LlamaGen assigns probability mass across the full codebook. VQGAN has a shorter curve because most nominal codes are inactive on ImageNet, and its active codes carry substantially higher average frequency.

\begin{figure}[t]
    \centering
    \includegraphics[width=0.78\textwidth]{ch04_fig_code_usage_sorted.png}
    \caption{Sorted ImageNet code-usage distribution. Codes are ranked by empirical frequency. The VQGAN curve ends near rank 972 because the remaining nominal codes are inactive on ImageNet.}
    \label{fig:rq1_code_usage_sorted}
\end{figure}

The top-$k$ columns in Table~\ref{tab:rq1_codebook_usage} report the cumulative
mass of the 10, 100, and 500 most frequent codes. High top-$k$ mass means that a
large fraction of the token sequence comes from a small set of IDs.

\section{Baseline Summary}
\label{sec:rq1_joint_view}

The baseline has two parts. First, LlamaGen gives better clean ImageNet reconstructions under the reported sample-level and distribution-level metrics. Second, the two tokenizers use their nominally equal codebooks very differently. Table~\ref{tab:rq1_joint_summary} keeps a compact version of these two facts together.

\begin{table}[t]
    \centering
    \scriptsize
    \begin{tabular}{lrrrrrrrr}
        \textbf{Tokenizer} & \textbf{PSNR} & \textbf{SSIM} & \textbf{LPIPS} & \textbf{FID} & \textbf{Active} & \textbf{Perplexity} & \textbf{Top-100} & \textbf{Codes for 99\%} \\ \hline
        LlamaGen & 20.79 & 0.558 & 0.228 & 2.19 & 16{,}384 (100.0\%) & 16{,}179.19 & 0.012 & 16{,}144 \\
        VQGAN    & 19.65 & 0.489 & 0.286 & 5.06 & 972 (5.9\%)       & 910.46    & 0.165 & 946 \\
    \end{tabular}
    \caption{Joint ImageNet summary of reconstruction quality and codebook usage. The final column reports the number of most frequent codes required to cover 99\% of token assignments.}
    \label{tab:rq1_joint_summary}
\end{table}

These results set the ImageNet baseline for later chapters. As expected from the
newer tokenizer's stated design goals, LlamaGen reconstructs the validation set
more accurately and uses nearly its full codebook. VQGAN reconstructs less
accurately and leaves roughly 94\% of its nominal codes inactive. Such codebook
underuse is a documented failure mode in discrete representation learning
~\cite{lancucki2020robust,yu2021vector,mentzer2023finite}, and LlamaGen
explicitly targets it through its tokenizer design~\cite{sun2024llamagen}.
